\documentclass[11pt]{article}
\usepackage[a4paper,margin=1in]{geometry}
\usepackage{amsmath,amssymb,amsthm}
\usepackage[T1]{fontenc}
\usepackage{lmodern}
\usepackage[hidelinks]{hyperref}
\setlength{\emergencystretch}{3em}

\theoremstyle{plain}
\newtheorem{theorem}{Theorem}
\newtheorem{lemma}[theorem]{Lemma}
\theoremstyle{definition}
\newtheorem{definition}[theorem]{Definition}
\theoremstyle{remark}
\newtheorem{remark}[theorem]{Remark}

\title{A Counterexample to Conjecture 00000008869:\\
Naming $2$ Loses in Sylver Coinage}
\author{%
  Feiyu\thanks{Independent researcher.}
  \and
  Claude (Opus 5.5)\thanks{Anthropic.}%
}
\date{2026-10-03}

\begin{document}
\maketitle

\begin{abstract}
Conjecture 00000008869 claims that in Sylver coinage naming $2$ is the first player's
winning first move, and that in the prime version (only primes and $1$ may be named)
the winning first move is still exactly $2$. Both claims are false. After $2$, the
second player names $3$; every integer $n\ge2$ is then a sum of $2$s and $3$s, so the
first player can only name $1$, which loses. The refutation is checked in Lean~4
against Mathlib.
\end{abstract}

\section{The conjecture as stated}

The original text of conjecture 00000008869 reads:

\begin{quote}
Definition: Sylver coinage: two players alternately name positive integers, naming
$1$ loses, and one may not name a nonnegative integer linear combination of
previously named numbers (Conway's coinage partial-order game). Conjecture: Naming
$2$ is the first player's winning first move; certification of the winning strategy
reduces to finiteness of the branching of the Sylvester gap spectrum of the semigroup
generated by the named set --- every position tree has finite ordinal rank; and in
the prime version (only primes and $1$ nameable) the winning first move is still
exactly $2$, the proof difference being that Frobenius numbers of odd-prime
subsemigroups have no uniform upper bound.
\end{quote}

The Chinese version in \texttt{conjectures/00000008869.md} states the same three
claims.

\section{Reading of the statement}

A \emph{position} is the finite set $S$ of numbers named so far. A positive integer
$n$ is a \emph{legal move} from $S$ if it is not a nonnegative integer combination of
elements of $S$, i.e.\ $n\notin\langle S\rangle$, the additive monoid generated by
$S$. The player who names $1$ loses; since $1\notin\langle S\rangle$ unless $1$ has
already been named, a player always has the move $1$ available, and naming it amounts
to resigning. In the prime version only primes and $1$ may be named.

The conjecture is a conjunction of three claims. We refute the first and the third:
\begin{enumerate}
\item[(C1)] naming $2$ is a winning first move in Sylver coinage;
\item[(C3)] in the prime version, the winning first moves are exactly $2$; in
  particular naming $2$ is a winning first move.
\end{enumerate}
Naming $2$ is a winning first move if the first player, after naming $2$, can force a
win against every reply. We use the standard inductive definition of game outcomes
(Definition~\ref{def:outcome}). The refutation is robust: after $2$ the second player
has a reply after which the first player's only legal move is $1$, so naming $2$ loses
under any definition of a winning strategy.

\begin{definition}\label{def:outcome}
For a set of allowed numbers, the player to move from $S$ \emph{wins} if some allowed
legal move $n\ne1$ leads to a position $S\cup\{n\}$ from which the player to move
\emph{loses}. The player to move from $S$ \emph{loses} if every allowed legal move
$n\ne1$ leads to a position from which the player to move wins. Both notions are
defined inductively (least fixed points); in particular no position is both won and
lost for the player to move.
\end{definition}

\section{Result}

\begin{theorem}\label{thm:main}
In Sylver coinage, and also in its prime version, the player to move from the
position $\{2\}$ wins by naming $3$. Hence naming $2$ is not a winning first move, and
(C1) and (C3) fail. In particular Conjecture 00000008869 is false.
\end{theorem}

\section{Proof}

\begin{lemma}\label{lem:legal3}
$3$ is a legal move from $\{2\}$.
\end{lemma}

\begin{proof}
$\langle 2\rangle$ consists of the even numbers, and $3$ is odd.
\end{proof}

\begin{lemma}\label{lem:only1}
From $\{2,3\}$ the only legal move is $1$.
\end{lemma}

\begin{proof}
Let $n\ge2$. If $n=2k$ then $n=k\cdot2$; if $n=2k+1$ with $k\ge1$ then
$n=(k-1)\cdot2+3$. Either way $n\in\langle2,3\rangle$, so $n$ is not legal.
\end{proof}

\begin{proof}[Proof of Theorem~\ref{thm:main}]
By Lemma~\ref{lem:only1}, the player to move from $\{2,3\}$ has no legal move other
than $1$, so that player loses (vacuously, in the sense of
Definition~\ref{def:outcome}). By Lemma~\ref{lem:legal3}, the player to move from
$\{2\}$ can name $3$, which is allowed in both versions since $3$ is prime. That move
leads to $\{2,3\}$, so the player to move from $\{2\}$ wins. After the first player
names $2$, it is the second player who moves from $\{2\}$, so the second player wins.
Since no position is both won and lost, the position $\{2\}$ is not lost for the player
to move, and naming $2$ is not a winning first move.
\end{proof}

\section{Formalization}

\sloppy
The Lean~4 project in \texttt{lean/} (file \texttt{lean/Submission/Basic.lean})
formalizes the game and proves Theorem~\ref{thm:main} against Mathlib.
\begin{itemize}
  \item \texttt{Legal S n} is $0<n$ and $n\notin$ \texttt{AddSubmonoid.closure S},
    the additive submonoid of $\mathbb N$ generated by $S$.
  \item \texttt{Outcome allowed b S} is the inductive outcome of
    Definition~\ref{def:outcome}: \texttt{b = true} means the player to move wins,
    \texttt{b = false} that the player to move loses. The parameter
    \texttt{allowed} is \texttt{anyNumber} for the usual game and
    \texttt{primeOrOne} for the prime version. \texttt{not\_win\_and\_lose} proves
    that no position is both won and lost.
  \item \texttt{NamingTwoWins allowed} states that $2$ is allowed and legal in the
    empty position and that \texttt{Outcome allowed false \{2\}} holds.
    \texttt{Claim1} is (C1), and \texttt{Claim3} is (C3) in its exact form: for
    every $n$, naming $n$ is a winning first move in the prime version iff $n=2$.
  \item \texttt{legal\_three} and \texttt{legal\_after\_two\_three} are
    Lemmas~\ref{lem:legal3} and~\ref{lem:only1}; \texttt{second\_player\_wins} gives
    \texttt{Outcome allowed true \{2\}} whenever $3$ is allowed.
  \item \texttt{not\_claim1} and \texttt{not\_claim3} are the refutations, and
    \texttt{conjecture\_00000008869\_false} states
    $\lnot(\texttt{Claim1}\wedge P\wedge\texttt{Claim3})$ for an arbitrary
    proposition $P$ standing for the second claim.
\end{itemize}
The toolchain is \texttt{leanprover/lean4:v4.33.1}, Mathlib is pinned to
\texttt{v4.33.1}, and \texttt{lake build} succeeds. The project contains no
\texttt{sorry}, no \texttt{admit} and no \texttt{native\_decide}, and introduces no
axioms: \texttt{\#print axioms} reports exactly \texttt{propext},
\texttt{Classical.choice} and \texttt{Quot.sound} for
\texttt{conjecture\_00000008869\_false} and \texttt{not\_claim3}.

\fussy
\section{Remarks}

\begin{remark}
By the same argument, naming $3$ also loses (the reply is $2$). By Hutchings'
theorem, naming any prime $p\ge5$ is a winning first move~\cite{ww}; our refutation
does not use this.
\end{remark}

\begin{remark}
We do not address the second claim (finite ordinal rank of position trees).
\end{remark}

\begin{thebibliography}{9}
\bibitem{ww} E.~R.~Berlekamp, J.~H.~Conway and R.~K.~Guy, \emph{Winning Ways for
  your Mathematical Plays}, 2nd ed., A K Peters, 2001--2004.
\end{thebibliography}

\end{document}
